\begin{textsample}{POS Dim 2 – vem\_ed\_33 – Score 2.00 – vem\_ed\_33/t178.txt}  \label{ex:f2_pos_147}
Novos projetos com Portugal Em 2025, a equipe de Apoio à Internacionalização do Ensino Superior (AIES) da Divisão de Extensão e Pesquisa no Ensino Superior (Depes) da Coordenadoria de Ensino Superior de Graduação (CGESG) do Centro Paula Souza iniciou novos Projetos Colaborativos Internacionais (PCIs) com duas instituições portuguesas: Instituto Politécnico de Viseu e Instituto Superior Politécnico de Gaya (ISP Gaya).
Além disso, \textbf{manteve} as colaborações com a Universidade de Aveiro, realizadas desde 2021 com as Fatecs Ipiranga e Lins. (Re)branding pequenos negócios: São Paulo e Viseu No primeiro semestre de 2025, as professoras Patrícia Patrício (Fatec Ipiranga) e Paula Pinheiro (Instituto Politécnico de Viseu) realizaram o PCI “(Re)branding pequenos negócios: São Paulo e Viseu”.
Grupos mistos de estudantes das duas instituições compararam estabelecimentos nas duas \textbf{cidades}: barbearias, cafés, confeitarias, pet shops, padarias, adegas e publicaram seus estudos no Padlet.
Os estudantes participaram de uma pesquisa sobre competências socioemocionais, realizada no início e no fim do projeto.
Os principais achados da pesquisa foram apresentados pelas professoras no Colóquio de Educação Internacional para o Sul Global (CEISG 2025), realizado em setembro — leia sobre o evento na edição 31 de VEm.
Um manual de negociação luso-brasileiro Diferenças entre estilos de negociação no Brasil e em Portugal foram identificadas pelos alunos participantes do PCI orientado pelas professoras Suillan Miguez Gonzalez, da Fatec Guaratinguetá, Paula Pinheiro e Isabel Rodrigues, do ESTGV/ Instituto Politécnico de Viseu.
A equipe de AIES / Depes / CGESG auxiliou no projeto, desde as reuniões iniciais com as professoras parceiras até o desenvolvimento do projeto, segundo a metodologia COIL (Colaborative Online International Learning).
Os grupos mistos de estudantes produziram um Manual de Negociação Luso-Brasileiro, mostrando algumas dessas diferenças de estilos: Clareza, objetividade e cumprimento dos compromissos Proximidade, confiança e abertura Maior peso da formalidade Informalidade facilita informação, mas pode ser vista como falta de rigor Confiança construída ao longo do tempo, com base em consistência e cumprimento de acordos Empresas brasileiras costumam começar com base em reuniões informais e confiança pessoal, ajustando detalhes depois

% matched lemmas: cidade, manter
\end{textsample}
